\section{The exact intersection of the radical and CM fields}
\label{sec:radical-cm-entanglement}

Put $k=\mathbb Q^{\mathrm{ab}}=\mathbb Q(\mu_\infty)$ and
\[
 K_3=\mathbb Q(\sqrt{-3}),\qquad K_4=\mathbb Q(i),\qquad
 F_{12}=K_3K_4=\mathbb Q(\zeta_{12}).
\]
Let $\mathcal P$ be the set of positive rational primes, and set
\[
 L_{\mathrm{rad}}=k(q^{1/n}:q\in\mathcal P,\ n\geq1),\qquad
 L_{\mathrm{CM}}=K_3^{\mathrm{ab}}K_4^{\mathrm{ab}}.
\]
All unqualified radicals are the positive real ones.  The presence
of all roots of unity makes these fields independent of that choice.

\begin{theorem}\label{thm:radical-cm-intersection}
One has
\begin{align}
 L_{\mathrm{rad}}\cap L_{\mathrm{CM}}
 &=L_{\mathrm{rad}}\cap F_{12}^{\mathrm{ab}}\notag\\
 &=k\bigl(q^{1/12}:q\in\mathcal P,\,
                         2^{1/8},3^{1/8}\bigr)=:I.
 \label{eq:radical-cm-intersection}
\end{align}
In particular
\begin{equation}
 \operatorname{Gal}(I/k)\simeq
 (\mathbb Z/12)^2\times
       \prod_{q\in\mathcal P\setminus\{2,3\}}\mathbb Z/6.
 \label{eq:radical-cm-common-group}
\end{equation}
The first two factors are indexed by $2,3$.  In these coordinates
conjugation by an element of $G_{\mathbb Q}$ is multiplication by
its cyclotomic character, reduced modulo the indicated orders.
\end{theorem}

\begin{proof}
We first compute the entire maximal subextension of $L_{\mathrm{rad}}$
that is abelian over $F_{12}$.  Recall the simultaneous radical
description
\[
 A_{\mathrm{rad}}=\operatorname{Gal}(L_{\mathrm{rad}}/\mathbb Q)
 =
 \left\{(u,(t_q)_q):
 \begin{array}{l}
 u\in\widehat{\mathbb Z}^{\times},\
 t_q\in\widehat{\mathbb Z},\\
 t_q\bmod2=\epsilon_q(u)\quad(q\in\mathcal P)
 \end{array}\right\},
\]
with multiplication $(u,t)(v,s)=(uv,t+us)$.
Here $\epsilon_q$ is the quadratic character of $\mathbb Q(\sqrt q)$,
written additively.  The translation subgroup over $k$ is
$T=\prod_q2\widehat{\mathbb Z}$.
This is the classical radical Galois description
\cite[\S2, Theorem~2.5]{MoreeStevenhagen}; the passage from finite
sets of primes to this product follows by compactness.
It also follows directly from the simultaneous Kummer-image lemma:
a positive rational radical lying in an abelian extension of
$\mathbb Q$ has its square in $\mathbb Q$, which is exactly the
displayed parity condition.

Write
\[
 U_{12}=\{u\in\widehat{\mathbb Z}^{\times}:u\equiv1\pmod{12}\},
 \qquad A_{12}=\operatorname{Gal}(L_{\mathrm{rad}}/F_{12}).
\]
Thus $A_{12}$ is the inverse image of $U_{12}$ in $A_{\mathrm{rad}}$.
Let $C_{12}=\overline{[A_{12},A_{12}]}$.  Since the multiplier quotient is
abelian, $C_{12}\subset T$.  For the convention
$[a,b]=aba^{-1}b^{-1}$, direct multiplication gives
\begin{equation}
 [(u,t),(v,s)]=
 \bigl(1,(1-v)t+(u-1)s\bigr).
 \label{eq:radical-cm-affine-commutator}
\end{equation}
We claim the exact simultaneous equality
\begin{equation}
 C_{12}=(24\widehat{\mathbb Z})_{\{2,3\}}\times
             \prod_{q\ne2,3}12\widehat{\mathbb Z}.
 \label{eq:radical-cm-derived-subgroup}
\end{equation}
The subscripts $\{2,3\}$ here index the rational-prime coordinates,
not their local components.

We make the distinction between these coordinates and the local
prime $\ell$ explicit in the proof.  Every closed subgroup of the
profinite abelian group $T$ is stable under
$\widehat{\mathbb Z}$-multiplication: approximate an element of
$\widehat{\mathbb Z}$ by integers and use closedness.
The idempotents of $\widehat{\mathbb Z}=\prod_\ell\mathbb Z_\ell$
therefore project $C_{12}$ onto its local components inside $C_{12}$ itself.
Consequently it suffices to determine each simultaneous local
subgroup $C_{12,\ell}$; there can be no additional entanglement between
distinct local primes.

Commuting a lift of $u\in U_{12}$ with a translation gives
$(u-1)T$.  At $\ell\geq5$ one may choose $u_\ell=2$ and all
other local components one.  Hence
$C_{12,\ell}=\prod_q\mathbb Z_\ell$.
At $\ell=3$, choosing $u_3=4$ gives
$\prod_q3\mathbb Z_3\subset C_{12,3}$.
Both multipliers in \eqref{eq:radical-cm-affine-commutator} are one
modulo three, proving the reverse inclusion.  Thus
$C_{12,3}=\prod_q3\mathbb Z_3$.
At $\ell=2$, choosing $u_2=5$ and using translations in
$\prod_q2\mathbb Z_2$ gives
\[
 \prod_q8\mathbb Z_2\subset C_{12,2}.
\]

It remains to determine the image modulo eight simultaneously
in every coordinate.  For $u\in U_{12}$ put
$a(u)=(u_2-1)/4\bmod2$.  Since $u_2\equiv1\pmod4$,
\[
 a(u)=\epsilon_2(u).
\]
Also $\epsilon_3(u)=0$, since $\sqrt3\in F_{12}$.
Reducing \eqref{eq:radical-cm-affine-commutator} modulo eight gives
in coordinate $q$
\begin{equation}
 4\bigl(\epsilon_2(u)\epsilon_q(v)
               -\epsilon_2(v)\epsilon_q(u)\bigr).
 \label{eq:radical-cm-mod-eight}
\end{equation}
It is zero for $q=2,3$, and is a multiple of four for every
other $q$.  This proves the necessary simultaneous upper bound.

For the reverse inclusion, the characters
$\epsilon_q|_{U_{12}}$ with $q\ne3$ are jointly independent.
Indeed a finite product would be trivial on $U_{12}$ precisely if
$\mathbb Q(\sqrt a)\subset F_{12}$, where $a$ is the corresponding
positive squarefree product of primes.  The three quadratic
subfields of $F_{12}$ are
$\mathbb Q(i)$, $\mathbb Q(\sqrt{-3})$, and $\mathbb Q(\sqrt3)$;
among positive squareclasses only $a=3$ occurs.  Excluding $3$
therefore gives independence.  Every finite evaluation map onto
the corresponding product of $\mathbb F_2$ is surjective.
Compactness then gives surjectivity onto the entire product.

Fix $u\in U_{12}$ with $\epsilon_2(u)=1$.  For an arbitrary family
$(b_q)_{q\ne2,3}\in\prod_{q\ne2,3}\mathbb F_2$, choose
$v\in U_{12}$ with
\[
 \epsilon_2(v)=0,\qquad \epsilon_q(v)=b_q\quad(q\ne2,3).
\]
Choose any lifts of $u,v$ to $A_{12}$.  Formula
\eqref{eq:radical-cm-mod-eight} now gives exactly the family
$4b_q$ in all the other coordinates, with zero at $2,3$.
This proves surjectivity onto the whole required product modulo
eight, not merely onto its individual projections.  Together
with $\prod_q8\mathbb Z_2\subset C_{12,2}$, it proves
\[
 C_{12,2}=(8\mathbb Z_2)_{\{2,3\}}\times
                       \prod_{q\ne2,3}4\mathbb Z_2.
\]
Combining the local equalities proves
\eqref{eq:radical-cm-derived-subgroup}.

The fixed field of $C_{12}$ is $L_{\mathrm{rad}}\cap F_{12}^{\mathrm{ab}}$.
In the radical coordinates, the subgroup
\eqref{eq:radical-cm-derived-subgroup} fixes precisely the
$24$th roots of $2,3$ and the $12$th roots of all other primes.
Since a $24$th root is generated by its eighth and twelfth powers
as root orders, this is the field $I$ in
\eqref{eq:radical-cm-intersection}.  More explicitly,
$q^{1/24}=(q^{1/8})/(q^{1/12})$.
The quotient of $2\widehat{\mathbb Z}$ by
$24\widehat{\mathbb Z}$ is $\mathbb Z/12$, and by
$12\widehat{\mathbb Z}$ is $\mathbb Z/6$.  This proves
\eqref{eq:radical-cm-common-group} and its conjugation action.

Each $K_j^{\mathrm{ab}}/F_{12}$ is abelian, so
$L_{\mathrm{CM}}\subset F_{12}^{\mathrm{ab}}$.  We have therefore
proved $L_{\mathrm{rad}}\cap L_{\mathrm{CM}}\subset I$.
For the reverse inclusion, all rational cube roots lie in
$K_3^{\mathrm{ab}}$, since $K_3$ contains $\mu_3$.
All rational fourth roots lie in $K_4^{\mathrm{ab}}$, since
$K_4$ contains $\mu_4$.

Let $\alpha=2^{1/8}$.  The field $\mathbb Q(i,\alpha)$ contains
\[
 \zeta_8=\frac{1+i}{\alpha^4},
\]
so it is the splitting field of $X^8-2$.  An automorphism fixing
$i$ sends $\alpha$ to $\zeta_8^a\alpha$ and sends $\zeta_8$ to
$\zeta_8^{1+4a}$.  Composition on the index $a\bmod8$ is
$a*b=a+b+4ab$, which is commutative.  Thus this extension is
abelian over $K_4$, and $2^{1/8}\in K_4^{\mathrm{ab}}$.

For the remaining radical put
$v=(2+\sqrt3)^{1/4}$ and $M=\mathbb Q(i,v)$.
The field $M$ contains $F_{12}$ and all conjugates
$i^a v^{\pm1}$, so it is Galois over $\mathbb Q$.
Every automorphism of $M$ fixing $\sqrt{-3}=i\sqrt3$ has the form
\[
 i\longmapsto i^\epsilon,\qquad
 v\longmapsto i^a v^\epsilon,
 \qquad a\in\mathbb Z/4,\quad\epsilon\in\{1,-1\}.
\]
Indeed the sign on $\sqrt3=v^4-2$ must equal the sign on $i$.
The product of two such automorphisms has exponent
$a\eta+\epsilon b$ and sign $\epsilon\eta$; this rule is
symmetric in $(a,\epsilon)$ and $(b,\eta)$.
Therefore $M/K_3$ is abelian and $v\in K_3^{\mathrm{ab}}$.

On the CM elliptic curve $E_4:Y^2=X^3-X$, the three-division
polynomial is $3X^4-6X^2-1$.  Choose a three-torsion point $(x,y)$
with $x^2=1+2\sqrt3/3$.  The curve equation gives
\[
 y^4=x^2(x^2-1)^2
       =\frac49\sqrt3(2+\sqrt3).
\]
All its torsion coordinates are abelian over $K_4$, since the
endomorphism by $i$ is defined there and the Tate-module action
commutes with that rank-one CM action
\cite[Proposition~12.13, printed p.~119]{MilneMF110}.
Thus $y\in K_4^{\mathrm{ab}}$.
As $\sqrt{2/3}\in k$, the element
\[
 z=\frac{y}{\sqrt{2/3}\,v}\in L_{\mathrm{CM}}
\]
satisfies $z^4=\sqrt3$ and hence $z^8=3$.
All eighth roots of unity already lie in $k$, so
$3^{1/8}\in L_{\mathrm{CM}}$.
This supplies every generator of $I$ and completes the proof.
\end{proof}

\subsection{The remaining compatibility condition on GT}

The arithmetic intersection theorem by itself does not combine
the two separate GT normalizations.  Here is the exact additional
condition it leaves.

Suppose the separate arithmetic quotient identifications are fixed:
\[
 G/N_{\mathrm{rad}}\simeq
        A_{\mathrm{rad}}=\operatorname{Gal}(L_{\mathrm{rad}}/\mathbb Q),
 \qquad
 G/C_{\mathrm{CM}}\simeq
        A_{\mathrm{CM}}=\operatorname{Gal}(L_{\mathrm{CM}}/\mathbb Q),
\]
with both identifications induced by the inclusion
$H=G_{\mathbb Q}\subset G$.
Restricting these quotient maps to $I$ gives two homomorphisms
$a,b:G\to A_I=\operatorname{Gal}(I/\mathbb Q)$, which agree on
$H$ and induce the same multiplier.  Their difference
\begin{equation}
 d(g)=a(g)b(g)^{-1}\in C=\operatorname{Gal}(I/k)
 \label{eq:radical-cm-compatibility-defect}
\end{equation}
is a continuous crossed homomorphism:
\[
 d(gh)=d(g)+\lambda_gd(h),\qquad d|_H=0.
\]
The addition uses the coordinates
\eqref{eq:radical-cm-common-group}.
The formula follows by writing $a(g)=d(g)b(g)$ and using the
scalar cyclotomic conjugation on $C$.

\begin{proposition}\label{prop:radical-cm-joint-gate}
The following are equivalent:
\begin{enumerate}
\item $G=(N_{\mathrm{rad}}\cap C_{\mathrm{CM}})H$;
\item the joint radical and CM image of $G$ equals that of $H$;
\item $a=b$, or equivalently $d=0$.
\end{enumerate}
Thus the remaining compatibility has exponent dividing twelve
over $k$: it consists of one cubic coordinate for every rational
prime, one quadratic coordinate for each prime other than $2,3$,
and one coordinate modulo four for each of $2,3$.
\end{proposition}

\begin{proof}
The first two assertions are equivalent by taking the common
kernel of the two quotient maps.
The joint Galois image is exactly
$A_{\mathrm{rad}}\times_{A_I}A_{\mathrm{CM}}$,
by the intersection theorem and the Galois correspondence.
The joint image of $G$ contains it, since $H\subset G$.
It is contained in it precisely when its two restrictions to
$A_I$ coincide.  This proves the equivalences.  The final list
is the primary decomposition of
\eqref{eq:radical-cm-common-group}.
\end{proof}

For example, on $G_1=\ker\lambda$, write the coordinate of the
CM restriction to $I$ as $\eta_q(g)$ in $\mathbb Z/n_q$, where
$n_2=n_3=12$ and $n_q=6$ for $q\ne2,3$.
Normalize it arithmetically by
$\eta_q(\sigma)=\kappa_q(\sigma)/2\bmod n_q$ for
$\sigma\in G_{\mathbb Q^{\mathrm{ab}}}$.
Then the required coordinate identities are precisely
\[
 \frac{\rho_q(g)}2=\eta_q(g)\pmod{n_q}
 \qquad(q\in\mathcal P,\ g\in G_1).
\]
They compare the rational Kummer cocycles with finite quotients
of the CM frames.  Separate surjectivity of the two quotient maps
does not prove these identities.
